\documentclass[12pt,a4paper]{article}
\usepackage[UTF8]{ctex}
\usepackage{amsmath}
\usepackage{amssymb}
\usepackage{geometry}
\geometry{left=2.5cm, right=2.5cm, top=2.5cm, bottom=2.5cm}
\usepackage{graphicx}
\usepackage{float}
\usepackage{bm}
\usepackage{siunitx}

\title{电容-电压法（C-V法）测量半导体杂质分布讲义}
\date{}

\begin{document}

\maketitle

作 \(\ln p - (1000 / T)\) 曲线和 \(\ln \left(\frac{p \cdot n}{T^3}\right) - (1000 / T)\) 曲线, 并且利用 \(p \cdot n = AT^3 \exp \left(\frac{-E_k}{kT}\right)\) , 求出硅的禁带宽度 \(E_g\) .

\section*{参考文献}

[1] E. H. 普特来. 霍耳效应及有关现象. 上海: 上海科学技术出版社, 1964.

[2] F. J. Morin, J. P. Maita. Electrical Properties of silicon Containing Arsenic and Boron. Phys. Rev., 1954.96:28.

[3] 吴思诚, 王祖铨. 近代物理实验. 2 版. 北京: 北京大学出版社, 1994.

[4] 黄昆, 谢希德. 半导体物理学. 北京: 科学出版社, 1958.

[5] L. J. Van der Pauw, Philips Technical Review, 1958~1959.20, 220.

[6] 夏平畴. 永磁机构. 北京: 北京工业大学出版社, 2000.

\section*{9-2 用电容-电压法测量半导体中的杂质分布}

半导体器件设计与制造的核心问题是如何控制半导体内的杂质分布, 以满足半导体器件实际应用的需要. 因此, 正确地检测杂质浓度及其分布是制备合格半导体器件的必要条件. 所以杂质浓度及其分布的测量是半导体材料及器件研究中的基本测量之一. 利用四探针或霍耳效应逐次去层测量薄层电导或薄层霍耳电压, 可以获得杂质分布以及迁移率随杂质浓度的变化. 但是, 这种逐次去层测量技术比较繁琐, 而且对样品具有破坏性. 通过测量不同直流偏压下 \(\mathrm{pn}\) 结势垒电容, 可以在不破坏器件本身情况下比较迅速地测得单边突变 \(\mathrm{pn}\) 结轻掺杂一边的杂质浓度及其分布. 例如, 利用这种方法可以迅速地求出用扩散工艺或者离子注入工艺所引入的杂质浓度及其分布. 然而, 对于重掺杂一边的杂质分布则需要用其他方法测定.

本实验的目的是: 利用 \(C - V\) 法测量 \(\mathrm{p}^+ \mathrm{n}\) 结轻掺杂一边的杂质浓度及其分布; 同时要了解锁定放大器的基本原理, 学会正确地使用锁定放大器.

\subsection*{一、原理}

\subsubsection*{(一) \(\mathrm{pn}\) 结势垒电容与杂质浓度的关系}

当 \(\mathrm{pn}\) 结中一边的杂质浓度比另一边大得多, 例如 \(N_{\mathrm{A}} \gg N_{\mathrm{D}}\) (\(N_{\mathrm{A}}\) 是 \(\mathrm{p}\) 区受主杂质浓度, \(N_{\mathrm{D}}\) 是 \(\mathrm{n}\) 区施主杂质浓度), 则在 \(\mathrm{p}\) 区和 \(\mathrm{n}\) 区交界处杂质分布有一突变, 这样的 \(\mathrm{p}^+ \mathrm{n}\) 结称为单边突变结, 浅扩散结常用它作近似. 在这单边突变结中, \(x_{\mathrm{n}} \gg x_{\mathrm{p}}\) (\(x_{\mathrm{p}}\) 和 \(x_{\mathrm{n}}\) 分别是 \(\mathrm{p}\) 区和 \(\mathrm{n}\) 区的空间电荷区宽度), 总空间电荷区宽度 \(w = x_{\mathrm{n}} + x_{\mathrm{p}} \approx x_{\mathrm{n}}\), 显然总空间电荷区内电势的变化几乎都落在轻掺杂的 \(\mathrm{n}\) 区, 而重掺杂的一边 \(x_{\mathrm{p}}\) 很小, 几乎可以忽略. 这时只需单独解轻掺杂一边的泊松方程, 就可以求得 \(\mathrm{pn}\) 结的特性. 如当 \(\mathrm{pn}\) 结上外加一反偏压 \(V_{\mathrm{R}}\) 时, 则空间电荷区宽度和偏压的关系仅与轻掺杂浓度 \(N_{\mathrm{D}}\) 有关

\[
w = \left[\frac{2 \epsilon \epsilon_{0}}{q N_{\mathrm{D}}} (V_{\mathrm{D}} + V_{\mathrm{R}})\right]^{1 / 2} \tag{9-2-1}
\]

式中 \(V_{\mathrm{D}}\) 为无外偏压作用下结两端间的电势差,称为接触电势差; \(\epsilon\) 为半导体的相对介电常量,硅的 \(\epsilon\) 为 \(11.8\); \(\epsilon_{0}\) 为真空电容率, \(\epsilon_{0} = (36 \pi \times 10^{11})^{- 1} \mathrm{~F} / \mathrm{cm}\); \(q\) 为电子电荷.这时, \(\mathrm{pn}\) 结的每一边所存储的电荷与空间电荷区宽度 \(w\) 成正比,即

\[
Q A = q A N_{\mathrm{D}}w = A\left[2q\epsilon \epsilon_{0}\left(V_{\mathrm{D}} + V_{\mathrm{R}}\right)N_{\mathrm{D}}\right]^{1 / 2}
\]

式中 \(A\) 为 \(\mathrm{pn}\) 结的结面积, \(Q\) 为单位结面积的空间电荷.单位面积的 \(\mathrm{pn}\) 结势垒电容

\[
\frac{C}{A} = \frac{\mathrm{d}Q}{\mathrm{d}V_{\mathrm{R}}} = \left[\frac{q\epsilon \epsilon_{0}N_{\mathrm{D}}}{2\left(V_{\mathrm{D}} + V_{\mathrm{R}}\right)}\right]^{1 / 2} = \frac{\epsilon \epsilon_{0}}{w} \tag{9-2-2}
\]

上式表明,当偏压的改变量 \(\mathrm{d}V_{\mathrm{R}}\) 足够小时,空间电荷区的电荷改变量 \(\mathrm{d}Q\) 与 \(\mathrm{d}V_{\mathrm{R}}\) 成正比,其比值就是 \(\mathrm{pn}\) 结势垒电容 \(C\). 它与一个金属平行板电容器有相似之处,电容的大小正比于 \(\mathrm{pn}\) 结的结面积 \(A\) ,反比于空间电荷区的宽度 \(w\). 当 \(\mathrm{pn}\) 结上外加偏压增加 \(\mathrm{d}V_{\mathrm{R}}\) 时,空间电荷区宽度也将增加 \(\mathrm{d}w\) ,原来在这 \(\mathrm{d}w\) 层内的载流子(n区中的电子,p区中的空穴)将流走,形成充放电电流,使空间电荷量增加 \(\mathrm{d}Q\);反之,当 \(\mathrm{pn}\) 结上外加偏压减少 \(\mathrm{d}V_{\mathrm{R}}\) 时,通过充放电电流使载流子填充到 \(\mathrm{d}w\) 层内,分别中和 \(\mathrm{n}\) 边这一层的电离施主的正电荷和 \(\mathrm{p}\) 边这一层内的电离受主的负电荷,使空间电荷区减薄 \(\mathrm{d}w\),从而使空间电荷量减少 \(\mathrm{d}Q\). 微量的充放电电荷 \(\mathrm{d}Q\) 集中在空间电荷区两边的薄层 \(\mathrm{d}w\) 内,这两个薄层相当于相距 \(w\) 的两个极板,半导体本身构成了电容器的电介质.但是 \(\mathrm{pn}\) 结势垒电容与平行板电容器也有区别,主要在于空间电荷区宽度 \(w\) 与平行板电容器的板间距离不同,空间电荷区宽度 \(w\) 不是固定不变的,它随外加偏压变化而变化.因此, \(\mathrm{pn}\) 结势垒电容也不是一个常量,它随外加偏压变化而变化.正由于具有这种特点,通常也把 \(\mathrm{pn}\) 结势垒电容称作微分电容.式(9- 2- 2)可以改写为

\[
\frac{1}{C^{2}} = \frac{2}{A^{2}q\epsilon \epsilon_{0}N_{\mathrm{D}}}\left(V_{\mathrm{R}} + V_{\mathrm{D}}\right) \tag{9-2-3}
\]

作出 \((1 / C^{2}) - V_{\mathrm{R}}\) 的图线——直线,由此直线的斜率可以求出施主杂质浓度 \(N_{\mathrm{D}}\) ,将此直线外推到电压轴,还可以求得接触电势差 \(V_{\mathrm{D}}\).

对于一个未知杂质分布的 \(\mathrm{pn}\) 结来说,可以利用 \(\mathrm{pn}\) 结电容- 电压曲线描绘出轻掺杂一边的杂质分布.如图9- 2- 1所示(图略),一个 \(\mathrm{p}^{+}\mathrm{n}\) 结在 \(\mathrm{n}\) 边有一个任意的杂质分布,当空间电荷区宽度变化 \(\mathrm{d}w\) 时,相应的单位面积的空间电荷变化量为

\[
\mathrm{d}Q = qN(w)\mathrm{d}w \tag{9-2-4}
\]

其中 \(N(w)\) 是空间电荷区宽度 \(w\) 边界处的杂质浓度.增加的电荷 \(\mathrm{d}Q\) 将引起电场改变 \(\mathrm{d}E\) ,由泊松方程可以求得[1]

\[
\mathrm{d}E = \frac{\mathrm{d}Q}{\epsilon \epsilon_{0}}
\]

相应的电势改变量

\[
\mathrm{d}\left(V_{\mathrm{R}} + V_{\mathrm{D}}\right) = \mathrm{d}V_{\mathrm{R}} = w\mathrm{d}E
\]

由以上两式消去 \(\mathrm{d}E\) ,得

\[
\frac{\mathrm{d}Q}{\mathrm{d}V_{\mathrm{R}}} = \frac{\epsilon \epsilon_{0}}{w} = \frac{C}{A} \tag{9-2-5}
\]

将以上三式代入式(9- 2- 4),并稍加整理得

\[
N(w) = \frac{2}{q\epsilon\epsilon_{0}A^{2}}\frac{1}{\frac{\mathrm{d}(1 / C^{2})}{\mathrm{d}V_{\mathrm{R}}}}
\]

由上式可以看出,当我们测量pn结势垒电容 \(c\) 随其偏压变化以后,作出 \((1 / C^{2}) - V_{\mathrm{R}}\) 曲线,并求出不同偏压下的 \(\mathrm{d}(1 / C^{2}) / \mathrm{d}V_{\mathrm{R}}\) ,代入上式,就可以求得 \(N(w)\).上式也可以改写为

\[
N(w) = \frac{1}{q\epsilon\epsilon_{0}A^{2}}\frac{C^{3}}{\mathrm{d}C / \mathrm{d}V_{\mathrm{R}}} \tag{9-2-6}
\]

采用直接求 \(c - V\) 曲线在不同偏压下的斜率 \(\mathrm{d}C / \mathrm{d}V_{\mathrm{R}}\) ,并将该偏压下的pn结电容 \(c\) 一起代入上式,也可以求得 \(N(w)\).此外,利用式(9- 2- 5)可以通过电容确定不同偏压下所对应的pn结空间电荷区宽度 \(w,N(w)\) 就是距离pn结交界 \((x = 0)\) 为 \(w\) 深处的杂质浓度, \(N(w) - w\) 关系曲线就是所要确定的杂质分布情况.我们利用这种 \(c - V\) 法可以求出扩散法或离子注入法所引入的杂质分布.但须注意,pn结中如果还有大量的陷阱中心或界面态存在,如硅中掺有金,以上的分析则需要加以修正.

\subsubsection*{(二) 测量势垒电容实验原理}

pn 结势垒电容是一个随直流偏压变化的微分电容,因此测量势垒电容时需要在pn 结上偏置一定的反向直流偏压 \(V_{\mathrm{R}}\) ,同时在 \(V_{\mathrm{R}}\) 上再叠加一个微小的交变电压信号 \(v(t)\). 这交变电压信号是来自锁定放大器中与参考信号同步的讯号,它经过待测的pn结电容 \(C_{\mathrm{x}}\) 和一个与之相串接的已知电容 \(C_{0}\) ,连接成如图9- 2- 2所示的测量线路.当 \(C_{0} \gg C_{\mathrm{x}}\) 时,在 \(C_{0}\) 两端的电压为

\[
v_{\mathrm{i}} = \frac{v(t)}{\frac{1}{j\omega C_{\mathrm{x}}} + \frac{1}{j\omega C_{0}}} \cdot \frac{1}{j\omega C_{0}} \approx \frac{C_{\mathrm{x}}}{C_{0}} v(t) \tag{9-2-7}
\]

上式表明电容 \(C_{0}\) 上的交变电压 \(v_{\mathrm{i}}\) 与待测的pn结电容 \(C_{\mathrm{x}}\) 成正比,比例系数 \(v(t) / C_{0}\) 等于 \(v_{\mathrm{i}} / C_{\mathrm{x}}\) 它表示将单位待测电容转换为电压信号的灵敏度.能否通过增大交变电压 \(v(t)\) 或降低 \(C_{0}\) 来提高这个灵敏度呢?我们知道,公式(9- 2- 7)成立的前提是 \(C_{0} \gg C_{\mathrm{x}}\). 此外,由于pn结是一个微分电容,交变电压 \(v(t)\) 必须比pn结上的压降 \(V_{\mathrm{D}} + V_{\mathrm{R}}\) 要小得多,否则将给微分电容的测量带来一定的误差.因此,我们只有选取适当的交变电压信号,经过检波后的输出幅度如果与输入电平成正比,那么,把检波后的直流信号接到 \(X - Y\) 记录仪(计算机)的 \(y\) 轴输入端, \(\mathrm{pn}\) 结上的反向偏压接到记录仪的 \(x\) 轴输入端,通过调节反向偏压的大小,就可以测得 \(\mathrm{pn}\) 结的 \(C - V\) 曲线.

\subsubsection*{(三) 锁定放大器的基本原理}

锁定放大器(LIA- lock- in amplifier)是一种交流电压表,它能精确测定深埋在噪声之中周期重复信号的幅值及相位,这种抑制噪声的作用主要是通过相敏检波器(PSD- phase sensitive detector)及低通滤波器(LPF- low pass filter)实现的.本实验采用了美国MODEL128A型锁定放大器,在它前面加一保护电路,具体的电路将在下面说明.

\paragraph{1. 相敏检波器}

相敏检波器实际上是一个模拟乘法器,如图9- 2- 3所示,输出电压 \(v_{0}\) 是输入电压 \(v_{i}\) 与参考信号 \(v_{\mathrm{R}}\) 的乘积,如果 \(v_{i}\) 和 \(v_{\mathrm{R}}\) 是两个正弦波,即

\[
v_{i} = E_{i}\sin (\omega_{1}t + \phi_{1})
\]
\[
v_{\mathrm{R}} = E_{\mathrm{R}}\sin (\omega_{2}t + \phi_{2})
\]

则输出电压 \(v_{0}\) 为

\[
v_{0}=v_{i}\cdot v_{\mathrm{R}}=\frac{E_{i}E_{\mathrm{R}}}{2}\cos\big[(\omega_{1}-\omega_{2})t+(\phi_{1}-\phi_{2})\big]-\frac{E_{i}E_{\mathrm{R}}}{2}\cos\big[(\omega_{1}+\omega_{2})t+(\phi_{1}+\phi_{2})\big]
\]

式中等式右边第一项为差频分量,第二项为和频分量.对于和频分量,可以用低通滤波器加以滤除.

相乘电路有许多形式,如开关型、电流控制型等等.

\paragraph{2. 低通滤波器}

在讨论低通滤波器之前,先简单回顾一下放大器噪声的问题

测量中的噪声是我们所不希望的无规则的杂乱信号,它是被测信号的自然背景.对微弱信号的放大和传递,影响最大的是器件内部的噪声及热噪声.热噪声是由于载流子无规则运动所引起的;器件内部的噪声除热噪声外,还有散弹噪声、闪烁噪声(1/f噪声)等.散弹噪声是由于载流子浓度梯度涨落引起的.散弹噪声和热噪声仅与测量频带宽度有关,而与频率无关.在直流到红外频率范围内,噪声谱是均匀分布的,即各个频率的噪声电压的平方平均值都是相同的,这类噪声通常称为白噪声,我们可以通过压缩频带宽度来减少它.适当选取参考信号的频率,可使1/f噪声忽略不计.

早期常用窄带滤波器压缩频带,但是由于窄带滤波器中心频率 \(f_{0}\) 不稳定,而带宽与 \(f_{0}\) 有关,因此抑制噪声的能力有限.近20年来所研制的种种锁相放大器,则是通过相敏检波器进行频谱变换,正如式(9- 2- 8)所描述的,由原来以 \(f_{1} = \omega_{1} / 2\pi\) 为中心的频谱,变换成以差频及和频为中心的两个频谱,再用低通滤波器压缩带宽,使大量的宽带噪声得到有效的抑制.

滤波器抑制噪声的能力通常用"等效噪声带宽"这一参数来描述.所谓等效噪声带宽是假定有一个与这个滤波器等效的理想滤波器存在,它具有绝对陡峭的通带.如果我们以相同的宽带噪声输入,在输出端得到相同的噪声幅值(均方根值),于是将此理想滤波器无衰减的带宽定义为原滤波器的等效噪声带宽.

低通滤波器一般采用性能好、体积小而且重量轻的有源滤波器,如图9- 2- 4所示.其输出与输入的关系为

\[
\frac{\dot{u}_{\mathrm{sc}}}{\dot{u}_{\mathrm{st}}} = \frac{R}{R_{0}}\frac{1}{1 + \mathrm{j}\omega R C} = -K\frac{1}{1 + \mathrm{j}\omega R C}
\]

因此,等效噪声带宽 \(\Delta f_{\mathrm{N}}\) 为

\[
\Delta f_{\mathrm{N}} = \frac{1}{K^{2}}\int_{0}^{\infty}\left|\frac{K}{1 + \mathrm{j}\omega R C}\right|^{2}\mathrm{d}f = \frac{1}{4R C}
\]

低通滤波器时间常量 \(T = R C,T\) 越长则等效噪声带宽 \(\Delta f_{\mathrm{N}}\) 越窄.但实际上由于漂移等问题, \(T\) 不能很长,于是 \(\Delta f_{\mathrm{N}}\) 也不可能无限小.

对于给定的白噪声,噪声电压正比于等效噪声带宽的平方根.因此,锁定放大器获得的信噪比改善,是输入信号的噪声带宽与低通滤波器噪声带宽之比的平方根,即

\[
\frac{\mathrm{输出信噪比}}{\mathrm{输入信噪比}} = \frac{\sqrt{\Delta f_{\mathrm{输入}}}}{\sqrt{\Delta f_{\mathrm{输出}}}}
\]

例如 \(\Delta f_{\mathrm{输入}} = 10 \mathrm{kHz}, T = 1 \mathrm{s}\) ,则 \(\Delta f_{\mathrm{输出}} = 1 / 4 T = 0.25 \mathrm{Hz}\) ,信噪比改善了200倍(46dB).如果输入的信噪比为1:10,则输出的信噪比将是20:1,一个深埋在噪声之中的微弱信号将被检测得到.

\paragraph{3. 相关器}

由于相敏检波器是一个模拟乘法器,低通滤波器实际上是一个积分器.因此,通过相敏检波及低通滤波器的输出信号 \(V_{0}\) 是输入信号与参考信号的乘积,再对时间积分,并取平均值有

\[
V_{0}^{\prime} = \lim_{T\to \infty}\frac{1}{2T}\int_{-T}^{T}V_{1}(t)\cdot V_{R}(t - \tau)\mathrm{d}t \tag{9-2-9}
\]

式中 \(\tau = (\phi_{1} - \phi_{2}) / \omega\) ,它是参考信号相对于输入信号的延迟时间,积分时间上限即低通滤波器的时间常量 \(T = R C,T\) 必须比参考信号的周期长得多.通常把式(9- 2- 9)所表示的 \(V_{0}\) 称为 \(V_{1}(t)\) 和 \(V_{\mathrm{R}}(t)\) 的相关函数,于是相敏检波器与低通滤波器的整体常统称为相关器或相关接收器.

考虑输入信号是一个埋没在噪声之中的微弱信号,即

\[
\left\{ \begin{array}{l}V_{\mathrm{S}} = E_{\mathrm{t}}\cos (\omega_{0}t + \theta) + E_{\mathrm{n}}\cos (\omega t + \alpha) \\ V_{\mathrm{R}} = E_{\mathrm{R}}\cos \omega_{0}t \end{array} \right. \tag{9-2-10}
\]

式中 \(\omega_{0}\) 为信号频率, \(\omega\) 为随机的噪声频率, \(\theta\) 为信号与参考信号的相位差, \(\alpha\) 为随机的噪声与参考信号的相位差.通过相敏检波器的输出电压

\[
\begin{array}{l}{{V_{0}=V_{\mathrm{S}}\cdot V_{\mathrm{R}}=\frac{1}{2}E_{\mathrm{t}}E_{\mathrm{R}}\cos\theta+\frac{1}{2}E_{\mathrm{t}}E_{\mathrm{R}}\cos(2\omega_{0}t+\theta)}}\\ {{+\frac{1}{2}E_{\mathrm{t}}E_{\mathrm{R}}\cos[(\omega-\omega_{0})t+\alpha]+\frac{1}{2}E_{\mathrm{t}}E_{\mathrm{R}}\cos[ (\omega+\omega_{0})t+\alpha]}}\end{array} \tag{9-2-11}
\]

通过低通滤波器的输出电压

\[
V_{0}^{\prime} = \frac{1}{2}E_{\mathrm{t}}E_{\mathrm{R}}\cos\theta + \frac{1}{2}E_{\mathrm{t}}E_{\mathrm{R}}\cos(\Delta \omega t + \alpha)
\]

式中 \(\Delta \omega = \omega - \omega_{0}\). 上式表明通过低通滤波器后, 式 (9- 2- 11) 所示的相敏检波器输出电压 \(V_{0}\) 中第二项及第四项, 都被低通滤波器滤除; 第一项为直流输出; 第三项中的 \(\Delta \omega\) 如果大于低通滤波器的带宽, 则也被滤除. 但是, 如果 \(\Delta \omega\) 小于低通滤波器的带宽, 则那部分噪声对输出 \(V_{0}\) 仍有影响, 如式 (9- 2- 12) 中等号后第二项所示.

实际的开关型相敏检波电路中, 参考信号为方波, 因为方波的幅度可以比较稳定, 当脉冲宽度为方波周期的一半, 并设参考信号的幅度 \(E_{\mathrm{R}} = 1\) , 用傅里叶级数展开:

\[
V_{\mathrm{R}} = \frac{4}{\pi}\sum_{n = 0}^{\infty}\frac{1}{2n + 1}\sin \left[(2n + 1)(2\pi f_{2}t + \phi_{2})\right] \tag{9-2-13}
\]

其中 \(n\) 为谐波数 \(f_{2}\) 为方波频率. 如果输入信号电压 \(V_{1}\) 仍为正弦波, 则相敏检波器的输出电压

\[
V_{0} = \sum_{n = 0}^{\infty}\frac{2E_{i}}{(2n + 1)\pi}\cos \left[2\pi [f_{1} - (2n + 1)f_{2}]t + [\phi_{1} - (2n + 1)\phi_{2}]\right] -\sum_{n = 0}^{\infty}\frac{2E_{i}}{(2n + 1)\pi}\cos \left[2\pi [f_{1} + (2n + 1)f_{2}]t + [\phi_{1} + (2n + 1)\phi_{2}]\right]
\]

上式表示输出电压 \(V_{0}\) 包括信号频率 \(f_{1}\) 与参考方波基频 \(f_{2}\) 的所有奇次谐波组成的和频与差频分量. 当 \(f_{1} = (2n + 1)f_{2}\) 时, 相敏检波器有同步输出, 也就是说, 开关型相敏检波器对任何一个奇次谐波都产生一个相敏直流输出, 其幅度与 \((2n + 1)\) 成反比. 设 \(f_{1}\) 与 \(f_{2}\) 的某一个奇次谐波同步, 如 \(f_{1} = (2n_{0} + 1)f_{2}\) , 则相关器的输出电压

\[
V_{0}^{\prime} = \frac{2E_{i}}{(2n_{0} + 1)\pi}\cos \left[\phi_{1} - (2n_{0} + 1)\phi_{2}\right] \tag{9-2-14}
\]

上式表明相关器输出的直流电位 \(V_{0}^{\prime}\) 与输入交流信号 \(E_{i}\) 成正比, 并与 \(\phi_{1} - (2n_{0} + 1)\phi_{2}\) 的余弦成正比. 当 \(\phi_{1} - (2n_{0} + 1)\phi_{2} = 0^{\circ}\) 或 \(180^{\circ}\) , \(V_{0}^{\prime}\) 正比于 \(E_{i}\) , 且为最大值; 当 \(\phi_{1} - (2n_{0} + 1)\phi_{2} = 90^{\circ}\) 或 \(270^{\circ}\) , \(V_{0}^{\prime}\) 为零. 图9- 2- 5画出了输入信号与参考信号处于不同相位的相敏检波器输出 \(V_{0}\) 的波形以及通过低通滤波器的输出 \(V_{0}^{\prime}\) 的电平. 由图9- 2- 5可以清楚地看出相敏检波器的相位检波特性.

如果输入信号为一幅值恒定并与参考方波频率相同 \((f_{1} = f_{2})\) 的方波信号, 则相关器输出电压

\[
V_{0}^{\prime} = \frac{2E_{i}}{\pi}\cos (\phi_{1} - \phi_{2}) = \frac{2E_{i}}{\pi}\cos \phi
\]

其中 \(E_{i}\) 为输入方波的幅度, \(\phi\) 为输入信号与参考信号之间的相位差. 上式也可以表达为

\[
V_{0}^{\prime} = \frac{2E_{i}}{\pi}\sin \left(\frac{\pi}{2} -\phi\right)
\]

如果相位差 \(\phi\) 接近 \(\frac{\pi}{2}\) , 上式可以写成

\[
V_{0}^{\prime} = E_{1}\left(1 - \frac{2\varphi}{\pi}\right) \tag{9-2-15}
\]

上式表明,输出电压 \(V_{0}^{\prime}\) 和信号与参考方波之间的相位差成线性关系,相关器可以作为鉴相器使用。式(9- 2- 14)和式(9- 2- 15)都说明相关器的输出含有输入信号的振幅和相位的信息,只要在参考信号的通道中有一个定标的移相器,就可以由输出信号确定输入信号的振幅和相位。

如果输入信号有噪声,与式(9- 2- 10)所示的参考信号为正弦波类似,在参考信号为方波的情况下,噪声与参考信号的差频所组成的各项中,只有落在低通滤波器的等效噪声带宽之内的那一小部分分量才对相关器的输出有影响。相关器的谐波响应如图9- 2- 6所示。相关器除了整流及相敏检波外,还起一个梳齿式滤波器的作用,即在所有 \(f_{2}\) 的奇次谐波频率的地方,都有一个带宽为 \(\Delta f_{\mathrm{N}} = 1 / 4RC\) 的带通滤波器,各带通滤波器输出的直流分量的相对幅度和方波的频谱特性一样。因此,相关器也是一个以参考信号频率为参数的方波匹配滤波器。若有淹没在噪声之中,并与参考信号相干的正弦波或方波信号输入,信号可以从相关器全部通过,噪声的绝大部分可以被滤除,这样就可以通过相关器把淹没在噪声之中的信号提取出来。

\paragraph{4. 同步积分器}

同步积分器的原理线路如图9- 2- 7所示.输入信号为正弦波信号

\[
I_{i} = \frac{E_{i}}{R}\sin (\omega t + \phi)
\]

图中S代表一个电子开关,它由参考方波 \(V_{\mathrm{R}}\) 所激励,参考方波的频率为 \(f_{\mathrm{R}}\) ,假设其幅度 \(E_{\mathrm{R}} = 1,V_{\mathrm{R}}\) 的表达式见式(9- 2- 13).两个电容 \(c\) 分别由电子开关S交替而又同步地与R相接,对输入信号进行积分.所谓同步是指参考信号与输入信号同步,它们的频率相等,同步积分器的输出电压[3] \(V_{0} = V_{01}x,x\) 为单位开关函数,其数学表示式为

\[
x = \frac{4}{\pi}\sum_{n = 0}^{\infty}\frac{1}{2n + 1}\sin (2n + 1)(2\pi f_{n}t)
\]

\(V_{01}\) 为同步积分器的输出方波的振幅.同步积分器同样可以看作是以参考信号频率为参量的方波匹配滤波器,具有很强的抗干扰能力和抑制噪声的能力,这里不再重复讨论.同步积分器与相关器的主要区别在于同步积分器的输出为交流信号,是一个与参考信号相同频率的方波,而相关器输出为一直流信号;另外,同步积分器的时间常数 \(T = 2R C\)

\paragraph{5. 锁定放大器的原理图}

此原理图如图9- 2- 8所示.放大器由信号通道、参考通道、同步积分器、相关器及电源等部分组成.同步积分器和相关器是该仪器的核心,起着抑制干扰和噪声的作用,并从噪声和干扰中检测信号.

相移器:要通过相关器的输出信号确定输入的待测信号的振幅和相位,必须改变已知参考信号的相位. 因此, 要求参考通道中有一个相移电路, 可在 \(0 \sim 360^{\circ}\) 之间连续调节.

图9- 2- 8电路中的各点波形如图9- 2- 9所示, 为了便于说明锁定放大器的原理, 这里假设了输入端伴随信号的干扰和噪声均不大. A处的正弦波是输入信号通过信号通道后的输出波形, 实际情况很可能是输入的正弦波信号淹没在干扰和噪声之中. B、D两处分别为参考通道输出的对称方波, 它们将分别去控制同步积分器和相关器的开关. C处是同步积分器输出的方波, E处为相乘电路的输出波形(全波整流波形), F处为低通滤波器输出的直流电压.

\subsection*{二、实验装置及测量线路}

\subsubsection*{(一)pn结 \(C - V\) 测量线路}

pn结 \(C - V\) 测量线路如图9- 2- 10所示.图中 \(C_{x}\) 为待测电容, \(C_{0}\) 有三个数值 \(500\mathrm{pF},5\times 10^{3}\mathrm{pF}\) 和 \(4.7\times 10^{4}\mathrm{pF}\) 可供选择.本实验中将测量一个平面扩散工艺制作的 \(\mathrm{p^{+}n}\) 结的势垒电容随偏压的变化.反向偏压为 \(10\mathrm{V}\) 时,反向电流小于 \(1\mu \mathrm{A}\) 在零偏压下,它的势垒电容约为几十皮法.假设 \(C_{x}\) 为 \(50\mathrm{pF}\) \(C_{0}\) 应远远大于 \(C_{x}\) ,因此,我们选取 \(C_{0}\) 为 \(5\times 10^{3}\mathrm{pF}\) 由于 \(C_{x}\) 的阻抗远远大于 \(C_{0}\) 的阻抗,故加在 \(C_{x}\) 和 \(C_{0}\) 上的交流电压基本上都降落在 \(C_{x}\) 上面.硅pn结上的压降 \(V_{0} + V_{\mathrm{th}}\) 大于 \(600\mathrm{mV}(V_{0}\) 为自建势),为了使 \(v(t)\ll V_{0} + V_{\mathrm{R}}\) ,选取 \(v(t)\) 为 \(40\mathrm{mV}\) 于是,根据式(9- 2- 7),在电容 \(C_{0}\) 两端的电压

\[
v_{\mathrm{i}}\approx \frac{C_{\mathrm{x}}}{C_{0}} v(t) = 400 \mu \mathrm{V} \tag{9-2-16}
\]

这也就是锁定放大器的输入信号, 它是一个幅度随 \(C_{x}\) 变化的交变信号. 当 \(\mathrm{p}^{+}\mathrm{n}\) 结上偏置的反向偏压变化时, \(C_{x}\) 将随偏压改变而变化. 锁定放大器输出的直流电压将与 \(C_{x}\) 成正比.

锁定放大器是一种相关测量, 需要外加上一个与输入信号相干的参考信号. 本实验中参考信号以及pn结上的交变信号 \(v(t)\) , 由锁定放大器内同步的参考信号提供.

锁定放大器的输入端加一保护电路,如图9- 2- 10所示.如果输入信号一旦大于 \(300\mathrm{mV}\) ,将使二极管 \(\mathbf{D}_{1}\) 或 \(\mathbf{D}_{2}\) 导通,以免仪器由于过载而损坏.

pn 结的直流偏置电路如图9- 2- 10所示, \(10\mathrm{~V~}\) 电源由一晶体管直流稳压电源供给,反向偏置电压可以通过连续调节电位器 \(R_{1}\) 而由 \(10\mathrm{~V~}\) 变化到 \(0.5\mathrm{~V~}\) 左右. \(100\mathrm{~k}\Omega\) 电阻 \(R_{2}\) 是为了防止交流信号短路而引入的(一般要求 \(R_{2} > > 1 / \omega C_{x}\) ),当pn 结反向电流小于 \(1\mu \mathrm{A}\) 时,由于没有直流通路,通过 \(100\mathrm{~k}\Omega\) 电阻的电压降可以忽略不计,则偏置电源相当于直接加在pn 结两端.

\subsubsection*{(二)使用锁定放大器的一些注意事项}

\begin{enumerate}
    \item 锁定放大器的输入是交流信号(正弦波,方波或类方波,即其脉冲宽度为信号周期的一半),输出是直流信号,输出与输入成正比,比例常量是锁定放大器的总增益.若被测信号不是交流信号,则需要调制成交流信号,例如利用斩波器进行斩波,使其变成类方波才能使用锁定放大器.
    \item 调制频率选择的原则,要根据传感器的时间响应、测量对象的信号频率以及锁定放大器的噪声性能来决定.例如,真空热电偶是一种传感器,响应时间为 \(20\mathrm{~ms}\) ,调制频率必须小于 \(50\mathrm{~Hz}\).在传感器时间响应许可的情况下,调制频率要选得大一些,尽量避免半导体器件的 \(1 / f\) 噪声及 \(50\mathrm{~Hz}\) 干扰.设被测信号最大有用频率为 \(f_{\mathrm{s max}}\) ,通常取调制信号频率 \(f_{\mathrm{R}} > 60f_{\mathrm{s max}}\).
    \item 锁定放大器的输入端往往是一个传感器,如光电倍增管、真空热电偶等,将微弱的非电量转换成弱电信号进行检测.因此,对传感器的性能,如灵敏度响应(即单位被测非电量的变化所引起的交变电压的大小)、响应时间、等效噪声功率以及输出阻抗等参量应有清楚的了解,从而适当地选取锁定放大器的放大倍数、调制频率以及前置放大器的输入阻抗等参量.
    \item 锁定放大器对噪声的抑制能力取决于积分器或低通滤波器的等效噪声带宽的大小.等效噪声带宽越小,抑制噪声的能力越强.等效噪声带宽又反比于积分器和低通滤波器的时间常量,因此锁定放大器信噪比的改善是以牺牲测量时间为代价的.

    对连续固定的信号测量,可采用长积分时间常量;对随时间变化的信号,则积分时间常量的设定必须与被测量的变化速度相适应.时间常量如果过大,虽然对减小噪声有利,但是有用信号中较高的频率部分被积分平滑而严重失真.设信号的最高频率为 \(f_{\mathrm{s max}}\) ,一般选取积分时间和低通时间常量 \(T \leqslant 1 / 3 f_{\mathrm{s max}}\).在讨论式(9- 2- 9)时曾指出,时间常量 \(T\) 必须比参考信号周期长得多,即时间常量既有上限,也有下限.究竟选取多大合适,应通过测量加以确定.

    \item 锁定放大器的两个最重要的指标是动态储备和输出漂移.它们是衡量锁定放大器的使用范围及工作性能的极限指标.我们知道,任何一个信号处理系统都包括下列三个临界电平:
    \begin{itemize}
        \item (1)最小可分辨信号电平(MDS);
        \item (2)满刻度信号输入电平(FS);
        \item (3)不相干信号的最大过载电平(OVL),即最大噪声输入.
    \end{itemize}
    由这三个电平确定输入总动态范围、动态储备以及输出的动态范围,如图9- 2- 11.

    动态储备是指在确定的灵敏度条件下,不相干信号允许的最大过载电平与满刻度输出所需的相干输入电平之比.它表示噪声或干扰比满刻度读数大多少分贝而锁定放大器还不过载, 它也决定了噪声或干扰比满刻度信号电平大多少倍时, 仪器还能从这些强干扰信号中检测出信号.

    输出动态范围是指在确定的灵敏度条件下, 给出满刻度输出时输入电平与最小能检出的相干输入电平之比. 它表示能测量的最小信号为满刻度读数的多少分之一.

    输入总动态范围是指在确定的灵敏度条件下, 允许的最大不相干电压与最小能检出的相干输入电压之比.

    由此可见, 如果用对数表示, 输入总动态范围等于动态储备与输出动态范围之和. 从设计锁定放大器角度考虑, 就是如何不断提高输入总动态范围的问题, 也就是既要提高动态储备, 又要减少直流漂移. 当直流放大器的元器件及电路选定以后, 漂移就由直流增益决定. 因此, 锁定放大器设计的中心问题是怎样既能得到足够大的交流增益, 又能得到大的动态储备. 从使用锁定放大器角度考虑, 输入总动态范围往往是给定的, 而满刻度信号输入电平的选择, 则根据所测对象的要求而定. 当被测信号噪声大时, 需适当减少交流增益, 即相敏检波器前的放大倍数需减少, 使有足够的动态储备, 避免噪声引起锁定放大器过载. 在总增益一定的情况下, 减少了交流增益, 就要增加相敏检波器后的直流增益, 即动态储备的提高是以增大漂移为代价的. 当被测信号的噪声不大时, 则可增加交流增益、减少直流增益、增加输出动态范围、相对地压缩动态储备而获得低漂移的正确测量. 有的锁定放大器设有高储备和高稳定的开关, 分别提供上述两种情况.

    \item 注意屏蔽及良好的接地
\end{enumerate}

\subsubsection*{(三)从噪声中提取微弱信号的线路}

图9- 2- 12所示的线路为一自制的噪声源, 它利用晶体管的噪声并进行放大. 噪声电压 \(v_{\mathrm{n}}\) 可达 \(500 \mathrm{mV}\) , 取其十分之一, 与一个 \(0.15 \mathrm{mV}\) 的正弦信号同时通过一个加法器叠加在一起, 再输入到锁定放大器, 见图9- 2- 13. 图中 \(v_{\mathrm{i}}\) 为 \(150 \mathrm{mV}\) , \(R_{2} / R_{11} = 1 / 10^{3}\) , \(R_{2} / R_{12} = 1 / 10\) . 锁定放大器的输入信号 \(v_{\mathrm{s}}\) 为:

\[
v_{\mathrm{s}} = -\left[\frac{R_{2}}{R_{12}} v_{\mathrm{n}} + \frac{R_{2}}{R_{11}} v_{\mathrm{i}}\right]
\]

先观测锁定放大器的输出电压 \(V_{0}\) , 然后去掉噪声信号, 再观测锁定放大器的输出电压. 这时将会发现, 先后两次的数值几乎相等. 这表明锁定放大器具有很强的抑制噪声能力. 一个深埋于噪声之中的微弱信号, 可以被锁定放大器检测得到.

\subsection*{三、实验内容}

\subsubsection*{(一)了解Model128Lock-inAmplifier的工作原理和使用方法}
\begin{enumerate}
    \item 阅读128Lock-inAmplifier简要说明书
    \item 给出安全使用128Lock-inAmplifier的方法和步骤
\end{enumerate}

\subsubsection*{(二)绘出以128Lock-inAmplifier为核心的 \(C - V\) 测试电路方框图}

\subsubsection*{(三)参照图9-2-10,测量信号 \(v(t)\) 的频率 \(f_{\mathrm{s}}\) 幅度 \(V_{\mathrm{so}}\) 和相位 \(\phi_{0}\)}
\begin{enumerate}
    \item 用示波器测量信号 \(v(t)\) 的频率 \(f_{\mathrm{s}}\) 和幅度 \(V_{\mathrm{so}}\)
    \item 用128Lock-inAmplifier测量信号 \(v(t)\) 的有效值 \(\overline{{V}}_{\mathrm{s}}\) 和相位 \(\phi_{0}\) .这里的 \(\phi_{0}\) 为选择开关K置"校 \(v(t)\) "位置时的测量结果
\end{enumerate}

\subsubsection*{(四)观测锁定放大器从噪声中提取微弱信号的作用}
\begin{enumerate}
    \item 用示波器观测白噪声,观测后把噪声幅度调至最小
    \item 用三通接头把白噪声和测量信号 \(v(t)\) 一同接到示波器,调白噪声使之完全淹没测量信号 \(v(t)\) ,将这种含噪声的 \(v(t)\) 信号接测试电路" \(N^{\prime \prime}\) 端(把选择开关K置"校 \(v(t)\) "位置).把锁定放大器的时间常数分别选为 \(1\mathrm{ms},100\mathrm{ms},1000\mathrm{ms}\) ,观测锁定放大器的输出信号的稳定性;增大白噪声,使锁定放大器输出不稳定,用 \(x - y\) 记录仪记录上述三个时间常量下的 \(V_{\mathrm{so}} - t\) 曲线( \(t\) 为时间).
\end{enumerate}

\subsubsection*{(五)检验锁定放大器输出( \(V_{0}\) )与被测电容( \(C_{x}\) )的线性关系,并观测相位 \((\phi_{x} - \phi_{0})\) 与 \(C_{x}\) 的关系.这里的 \(\phi_{x}\) 为选择开关K置"测量"位置时,逐一接入已知待测电容 \(C_{x}\) 时的测量结果.注意:在实验中校准(或测定) \(v(t)\) 时,务必先置锁定放大器灵敏度于 \(250\mathrm{mV}\) 挡!因为 \(v(t)\) 是 \(100\mathrm{mV}\) 左右的信号!}

\subsubsection*{(六)研究测试电路输出信号相位 \(\phi\) 与二极管反向偏压 \(V_{\mathrm{R}}\) 之关系}
\begin{enumerate}
    \item 把待测二极管( \(\mathbf{p}^{\star}\mathbf{n}\) 结)接入测试电路
    \item 把选择开关K置于"测量"位置,逐渐增大锁定放大器灵敏度至适当的灵敏度挡,测量二极管在不同反偏压(0至 \(9\mathrm{V}\) )下测试电路输出信号的相位 \(\phi_{v}\) 值.
    \item 照上述步骤,逐一测量 \(D_{1}\) 至 \(D_{5}\) 各二极管的相位 \(\phi_{v}\) ,观测相位 \((\phi_{v} - \phi_{0})\) 与偏压 \(V_{\mathrm{R}}\) 的关系(数据列成表格).
    \item 试说明各个二极管的相位 \((\phi_{\mathrm{v}} - \phi_{0}) - V_{\mathrm{R}}\) 关系为什么不一样?试考虑当二极管出现较大的反向漏电流时对 \((\phi_{\mathrm{v}} - \phi_{0})\) 的影响.什么样的 \(\mathrm{p}^{\star}\mathrm{n}\) 单边突变结二极管可用作 \(C - V\) 法测定掺杂分布?
\end{enumerate}

\subsubsection*{(七) \(C - V\) 测量}
\begin{enumerate}
    \item 用 \(X - Y\) 记录仪记录 \(C - V\) 曲线

    把选择开关K置于"测量"挡,把用作 \(C - V\) 法测定掺杂分布的 \(\mathrm{p}^{\star}\mathrm{n}\) 单边突变结二极管接入电路,适当选择锁定放大器的灵敏度和时间常量、 \(X - Y\) 记录仪的 \(x\) 轴和 \(y\) 轴零点、记录仪的灵敏度.缓慢地调节反偏压旋钮,使记录仪描出尽可能光滑的 \(C - V\) 曲线,直到偏压 \(V_{\mathrm{R}}\) 升至 \(9\mathrm{V}.C - V\) 曲线画毕以后,抬起记录笔,重置 \(V_{\mathrm{R}} = 0\)

    \item 用标准电容标定 \(C - V\) 曲线的 \(C\) 轴

    在二极管处换接上一个标准电容,记录下相应的锁定放大器输出幅度.

    \item 修正 \(C - V\) 曲线的零点

    本实验样品盒上标定 \(1.4\mathrm{pF}\) 的挡,这实际上为空挡(即没接任何元件,为开路状态),但毕竟具有一定的电容.在 \(C - V\) 曲线中应扣除这一小电容.
\end{enumerate}

\subsubsection*{(八)关机}
\begin{enumerate}
    \item 把128Lock-inAmplifier锁定放大器的相移置于 \(^{\ast}0^{\circ}\) 位置;灵敏度置 \(^{\ast}250\mathrm{mV}^{\ast}\) 挡;锁定放大器的输入置"B"通道.然后关闭锁定放大器电源.
    \item 关闭 \(X - Y\) 记录仪.
    \item 逐一关闭直流电源、PZ8直流数字电压表、示波器的电源
    \item 关所用仪器电源
\end{enumerate}

\subsection*{四、预习要求}

\begin{enumerate}
    \item 学习了解杂质浓度分布与 \(\mathrm{p}^{\star}\mathrm{n}\) 结势垒电容之间的关系
    \item 本实验是如何测定pn结势垒电容以及势垒电容随负偏压的变化?
    \item 了解锁定放大器从噪声中提取微弱信号的工作原理
    \item 掌握安全使用锁定放大器的要点
\end{enumerate}

\subsection*{五、思考题和实验报告要求}

\subsubsection*{(一)实验过程思考题}
\begin{enumerate}
    \item 如何正确地选择锁定放大器灵敏度?锁定放大器的放大倍数是如何确定的?怎样估计从锁定放大器输出信号的大小?如何测定信号的相位?
    \item 在使用锁定放大器时,对所提取微弱信号的频率与模拟乘法器中参考信号的频率有何要求?噪声中与参考信号频率相同的成分能否用锁定放大器来消除它的影响?
    \item 如果 \(\mathrm{p}^{\star}\mathrm{n}\) 单边突变结二极管存在漏电,对测试结果会有什么影响?
\end{enumerate}

\subsubsection*{(二)实验报告的要求及思考题}
\begin{enumerate}
    \item 简述锁定放大器从噪声中提取微弱信号的工作原理
    \item 概述利用 \(C - V\) 曲线测定单边突变 \(\mathrm{p}^{\star}\mathrm{n}\) 结轻掺杂一边的杂质浓度及其分布的原理
    \item 说明有的二极管相位 \((\phi_{\mathrm{v}} - \phi_{0})\) 随 \(V_{\mathrm{R}}\) 而变化的原因.
    \item 利用 \(C - V\) 曲线计算二极管中浅掺杂一边的杂质浓度分布.

    列表给出各偏压 \(V_{\mathrm{R}}\) 下的电容 \(C,1 / C^{2},\mathrm{d}C / \mathrm{d}V,w\) 及杂质浓度 \(N(w)\) 值(在本实验中已知 \(\mathbf{p}^{\star}\mathbf{n}\) 结面积 \(A = 5.03\times 10^{- 3}\mathrm{cm}^{2}\) \(\epsilon \epsilon_{0} = 11.8\times 8.854\times 10^{- 14}\mathrm{F / cm})\) ,并作出 \(N(w) - w\) 杂质分布曲线.

    \item 作 \(1 / C^{2} - V_{\mathrm{R}}\) 关系曲线,由此曲线的直线部分外推至 \(1 / C^{2} = 0\) ,外推直线在 \(V\) 轴之截距即为单边突变结的自建势 \(V_{\mathrm{D}}\) ,给出 \(V_{\mathrm{D}}\) 的数值.
    \item 按实验内容顺序给出原始实验数据,并做相应的说明.
\end{enumerate}

\section*{参考文献}

[1] 吴思诚,王祖铨.近代物理实验.2版.北京:北京大学出版社,1995.

[2] E. Schibli, A. G. Milnes. Effects of deep impurities on \(\mathbf{n} + \mathbf{p}\) junction reverse- biased small- signal capacitance. Solid State Electron, 1968,11:323.

[3] 陈佳圭.如何正确使用锁相放大器.物理,1982,11(5):287.

\end{document}